\documentclass[10pt, aspectratio=169]{beamer}

% --- Configuración del Tema ---
\usetheme{metropolis}
\usecolortheme{owl} % Fondo oscuro, alto contraste
\usepackage[spanish]{babel}
\usepackage[utf8]{inputenc}
\usepackage{amsfonts, amsmath, amssymb, amsthm}
\usepackage{graphicx}
\usepackage{tikz}

% --- Colores Personalizados ---
\definecolor{diagonalred}{RGB}{140, 20, 40}
\definecolor{mitred}{RGB}{163, 31, 52}

% --- Información del Título ---
\title{Los Límites de la Razón}
\subtitle{Sesión 3: Lógica de primer orden y teoría de modelos}
\author{Rosh Guadiana}
\date{Junio 2026}

\begin{document}

% --- Slide 1: Título ---
\maketitle

% --- Slide 2: El Puente Roto (Parte 1) ---
\begin{frame}{El Puente Roto: La insuficiencia de Boole}
    \begin{center}
        \Large \textit{"Todo número par mayor que 2 es la suma de dos primos"}
        
        \vspace{0.8cm}
        \textbf{Traducción en Lógica Proposicional:}
        
        \vspace{0.4cm}
        \Huge $p$
    \end{center}
\end{frame}

% --- Slide 3: El Puente Roto (Parte 2) ---
\begin{frame}{La necesidad de un bisturí}
    \vspace{0.5cm}
    \Large
    La Lógica Proposicional es opaca; trata pensamientos hipercomplejos como bloques indivisibles. 
    
    \vspace{0.8cm}
    \alert{\textbf{El siguiente paso:}} \\
    \normalsize
    \vspace{0.3cm}
    Para que la matemática se autoanalice, necesitamos romper el átomo proposicional usando \textbf{cuantificadores} ($\forall, \exists$) y \textbf{variables} para ver su estructura interna.
\end{frame}

% --- Slide 4: La Asfixia del Vacío (Parte 1) ---
\begin{frame}{La Asfixia del Vacío}
    \begin{center}
        \Large ¿Verdadero o Falso?
        
        \vspace{1.2cm}
        \huge
        $\forall x \exists y \ R(x, y) \to \exists y \forall x \ R(x, y)$
    \end{center}
\end{frame}

% --- Slide 5: La Asfixia del Vacío (Parte 2) ---
\begin{frame}{El pecado capital de la lógica}
    \vspace{0.5cm}
    \Large
    Inyectarle un "mundo" a símbolos tipográficos completamente vacíos.
    
    \vspace{0.8cm}
    \normalsize
    En la Lógica de Primer Orden (FOL), los cuantificadores y las variables no valen nada sin una \textbf{Firma Lógica} y un \textbf{Dominio}. 
    
    \vspace{0.4cm}
    \alert{El sistema formal evalúa diseños que tienen forma, pero no contenido.}
\end{frame}

% --- Slide 6: Sintaxis y Firma Lógica (Parte 1) ---
\begin{frame}{Sintaxis y la Firma Lógica ($\mathcal{L}$)}
    \begin{center}
        \huge $\mathcal{L} = (\mathcal{F}, \mathcal{P}, ar)$
    \end{center}
    
    \vspace{0.8cm}
    \Large \textbf{El vocabulario de la máquina:}
    \normalsize
    \vspace{0.3cm}
    \begin{itemize}
        \setlength\itemsep{0.3cm}
        \item $\mathcal{F}$: Conjunto de símbolos de \textbf{funciones}.
        \item $\mathcal{P}$: Conjunto de símbolos de \textbf{predicado}.
        \item $ar$: La \textbf{aridad} (número de argumentos) de cada símbolo.
    \end{itemize}
\end{frame}

% --- Slide 7: Sintaxis y Firma Lógica (Parte 2) ---
\begin{frame}{El estado de las variables}
    \vspace{0.5cm}
    \Large \alert{\textbf{Anatomía de las variables:}}
    
    \vspace{0.6cm}
    \normalsize
    \begin{itemize}
        \setlength\itemsep{0.6cm}
        \item \textbf{Variables Libres:}\\ 
        Esperando recibir un valor (un significado).
        \item \textbf{Variables Ligadas:}\\ 
        Atrapadas por la dictadura de un cuantificador ($\forall x$ o $\exists x$).
    \end{itemize}
\end{frame}

% --- Slide 8: Semántica (Parte 1) ---
\begin{frame}{La Semántica y el Mundo}
    \begin{center}
        \vspace{0.5cm}
        \Large \textbf{Estructura} $\mathfrak{A}$ = Dominio + Interpretación
        
        \vspace{1.2cm}
        \huge
        $\mathfrak{A} \models \varphi$
    \end{center}
\end{frame}

% --- Slide 9: Semántica (Parte 2) ---
\begin{frame}{El anclaje del significado}
    \vspace{0.5cm}
    \Large
    Para que la sintaxis cobre vida, necesitamos inyectarle un \textbf{Universo de Discurso} (Dominio) y mapear cada símbolo vacío a una relación real en ese universo.
    
    \vspace{0.8cm}
    \normalsize
    \alert{Las verdades matemáticas no flotan en el vacío; existen \textbf{únicamente} relativas a un modelo.}
\end{frame}

% --- Slide 10: El Programa de Hilbert (Parte 1) ---
\begin{frame}{El Delirio de Omnisciencia}
    \begin{columns}
        \begin{column}{0.4\textwidth}
            \begin{center}
                % Recuerda ajustar la ruta de esta imagen
                \includegraphics[height=0.6\textheight]{ruta/a/hilbert.jpg} \\
                \vspace{0.2cm}
                \textbf{David Hilbert}
            \end{center}
        \end{column}
        \begin{column}{0.6\textwidth}
            \begin{center}
                \Large \textit{"Wir müssen wissen. \\ Wir werden wissen."}
                
                \vspace{0.4cm}
                \normalsize (Debemos saber. Sabremos.)
            \end{center}
        \end{column}
    \end{columns}
\end{frame}

% --- Slide 11: El Programa de Hilbert (Parte 2) ---
\begin{frame}{El sueño absolutista (1920s)}
    \vspace{0.5cm}
    \Large \textbf{El Programa de Hilbert:}
    
    \vspace{0.6cm}
    \normalsize
    Mecanizar toda la matemática usando el lenguaje impecable de la Lógica de Primer Orden (FOL). 
    
    \vspace{0.6cm}
    \alert{\textbf{El objetivo:}} \\
    Erradicar cualquier intuición humana y probar de forma puramente sintáctica que el sistema estaba libre de contradicciones.
\end{frame}

% --- Slide 12: El Clímax (Parte 1) ---
\begin{frame}{La victoria efímera}
    \begin{center}
        \Large \textbf{Teorema de Completitud de Gödel (1929)}
        
        \vspace{0.8cm}
        \huge
        $\Gamma \models \varphi \iff \Gamma \vdash \varphi$
    \end{center}
    
    \vspace{0.8cm}
    \normalsize
    La máquina sintáctica ($\vdash$) logra capturar toda la semántica ($\models$). El mundo creyó que Hilbert había ganado.
\end{frame}

% --- Slide 13: El Clímax (Parte 2) ---
\begin{frame}{El Susurro que Mató a Hilbert}
    \vspace{0.5cm}
    \Large \textbf{Congreso de Königsberg (1930):}
    
    \vspace{0.6cm}
    \normalsize
    Mientras Hilbert daba un discurso glorioso proclamando la victoria de su programa, en una mesa redonda, un joven Gödel anunciaba en voz baja su \textit{Teorema de Incompletitud}... 
    
    \vspace{0.8cm}
    \alert{\Large El sueño de Hilbert estaba matemáticamente muerto.}
\end{frame}

\end{document}
